\begin{solapartat}
Del gràfic ha completat 7 voltes.

Del esquema la longitud del eix major és: $35R_S + 2R_S + 200R_S = 237R_S$

Si el resultat és $235R_S$ perquè no s'ha tingut en compte el diàmetre del Sol, llavors la qualificació de la resposta serà 0,6~p enlloc de 0,8~p.

\figura[0.5]{sol_fig1.png}
\end{solapartat}

\begin{solapartat}
Dibuix d'una òrbita el·líptica amb el Sol en un dels seu focus.

Identificar els punts A i C com el més proper i el més allunyat del Sol.

Representar correctament la direcció i sentit dels vectors acceleració i velocitat.

Conservació de l'energia mecànica:
% DUBTE: l'original escriu M_T (massa de la Terra) on hauria de ser la massa del Sol
\[ E_m = \frac{1}{2}mv^2 - G\,\frac{mM_T}{r} \Longrightarrow \frac{1}{2}mv^2 = \mathit{constant} + G\,\frac{mM_T}{r} \]
Quan $r$ és mínim, llavors el terme $G\frac{mM_T}{r}$ i l'energia cinètica són màxims, per tant, el punt on la velocitat és més gran és el punt més proper al Sol, el punt A. També es pot argumentar a partir de la conservació de l'energia mecànica que l'energia cinètica és màxima quan l'energia potencial gravitatòria és mínima, i això passa quan la distància entre el Sol i la nau és mínima.

\destaca{Alternativament}

Segona llei de Kepler: ``Un segment rectilini que uneix la nau i el Sol escombra àrees iguals en intervals de temps iguals''.

Com la longitud d'aquest segment és mínima al punt A, llavors per escombrar la mateixa àrea en el mateix interval de temps, cal que el desplaçament sigui màxim, és a dir, cal que la velocitat de la nau sigui màxima.

El moment angular de la sonda respecte el centre del Sol s'expressa com:
\[ \vec{L}_{Parker,Sol} = m\vec{r}\times\vec{v}, \]
on $\vec{r}$ és el vector de posició que apunta del centre del Sol a la sonda, $m$ és la massa de la sonda i $\vec{v}$ és la seva velocitat. El mòdul del moment angular de la sonda respecte el Sol, s'expressa com:
\[ L_{Parker,Sol} = m\,r\,v\sin\theta. \]
Com el moment angular i la massa són constants:
\[ v = \frac{\mathit{constant}}{r\sin\theta} \]
La velocitat serà màxima en el punt on el producte $r\sin\theta$ sigui mínim, i aquest punt és el periheli, atès que és el punt de la òrbita més proper al Sol.
\end{solapartat}
